% !TeX spellcheck = en_US
\documentclass[french]{yLectureNote}

\title{Introduction à la termodynamique}
\subtitle{Thermodynamique}
\author{Paulhenry Saux}
\date{\today}
\yLanguage{Français}

\professor{}

\usepackage{graphicx}%----pour mettre des images
\usepackage[utf8]{inputenc}%---encodage
\usepackage{geometry}%---pour modifier les tailles et mettre a4paper
%\usepackage{awesomebox}%---pour les boites d'exercices, de pbq et de croquis ---d\'esactiv\'e pour les TP de PC
\usepackage{tikz}%---pour deiffner + d\'ependance de chemfig
\usepackage{tkz-tab}
\usepackage{tabularx}%---pour dimensionner automatiquement les tableaux avec variable X
\usepackage{awesomebox}%---Pour les boites info, danger et autres
\usepackage{menukeys}%---Pour deiffner les touches de Calculatrice
\usepackage{fancyhdr}%---pour les en-t\^ete personnalis\'ees
\usepackage{blindtext}%---pour les liens
\usepackage{hyperref}%---pour les liens (\`a mettre en dernier)
\usepackage{caption}%---pour la francisation de la l\'egende table vers Tableau
\usepackage{pifont}
\usepackage{array}%---pour les tableaux
\usepackage{lipsum}
\usepackage{yFlatTable}
\usepackage{multicol}
\newcommand{\Lim}[1]{\lim\limits_{\substack{#1}}\:}
\renewcommand{\vec}{\overrightarrow}
\newcommand{\bz}{\overline{z}}
\DeclareMathOperator{\card}{card}
\renewcommand{\vec}{\overrightarrow}
\newcommand{\N}[0]{\mathbb{N}}
\newcommand{\R}[0]{\mathbb{R}}
\newcommand{\C}[0]{\mathbb{C}}
\newcommand{\dd}[0]{\mathrm{d}}
\newcommand{\er}{\vec{e_{\rho}}}
\newcommand{\ep}{\vec{e_{\varphi}}}
\newcommand{\pip}{\pi^+}
\newcommand{\pim}{\pi^-}
\newcommand{\mv}{\mathcal{V}}
\DeclareMathOperator{\Div}{Div}
\newcommand{\rot}{\vec{\mathrm{rot}}}
\newcommand{\grad}{\vec{\mathrm{grad}}}
\begin{document}
\setcounter{chapter}{3}
\chapter{Second principe de thermodynamique}

\section{Entropie}

\begin{theorem}
    Comme l'intégrale de $\frac{\delta Q_{rev}}{T}$ sur un cycle fermé est nulle, la quantité $\int_A^B \frac{\delta Q_{rev}}{T}$ est indépendante du chemin suivi. Elle définit la variation d'une fonction d'état appelée \textbf{Entropie}, notée $S$.
    $$\Delta S_{AB} = S_B - S_A = \int_A^B \frac{\delta Q_{rev}}{T}$$
\end{theorem}

\section{Énoncé du second principe}
\begin{theorem}[Énoncé détaillé]
    $$\Delta_{i\to f} = S_f-S_i = \int_i^f \frac{\delta Q_{QS}}{T} = -\Delta S_{ext, if}+S^p_{i\to f}$$
    avec 
    \begin{itemize}
        \item $-\Delta S_{ext,if} = 0$ si la transformation est adiabatique
        \item $-\Delta S_{ext,if} = \frac{Q_{i\to f}}{T_{source}}$ si la transformation s'effectue au contact d'une source.
    \end{itemize}
      et
      \begin{itemize}
        \item $S^p_{if}=0$ si c'est réversible
        \item $S^p_{if}>0$ si c'est irréversible
      \end{itemize}
\end{theorem}
\tipsInfo{Signification physique}{  $S_e = \int \frac{\delta Q}{T_{ext}}$ est l'entropie échangée avec le milieu extérieur.
        et $S_p$ est l'entropie produite (ou créée) au sein du système.
    }

\begin{theorem}[Entropie de l'Univers]
    Pour un système isolé (ou l'Univers entier), $Q=0 \implies S_e = 0$.
    Par conséquent : $\Delta S_{univers} \geq 0$.
    L'entropie d'un système isolé ne peut que croître.
\end{theorem}

\section{Énoncé de Thomson / de Kelvin}

\criticalInfo{Rappel : Moteur}{Pour un cycle moteur, le travail total échangé est négatif ($W < 0$), ce qui signifie que le système fournit de l'énergie mécanique au milieu extérieur.}

\begin{theorem}[Énoncé de Thomson]
    Un système en contact avec une seule source de chaleur (cycle monotherme) ne peut pas fournir de travail au milieu extérieur au cours d'un cycle. 
    Autrement dit : pour un cycle monotherme, $W \geq 0$.
\end{theorem}

\begin{theorem}[Conséquence]
    Il est impossible de transformer intégralement de la chaleur en travail de façon cyclique. En revanche, la transformation intégrale de travail en chaleur est possible (ex: effet Joule).
\end{theorem}

\section{Exemples d'application}
\subsection{Calcul de $\Delta S$ entre deux états quelconques}

L'entropie étant une fonction d'état, on peut choisir n'importe quel chemin fictif quasi-statique entre $(P_0, V_0, T_0)$ et $(P_1, V_1, T_1)$.

\begin{proposition}[Formules de l'entropie du Gaz Parfait]
    Selon les variables d'état choisies, on utilise les identités thermodynamiques :
    \begin{itemize}
        \item \textbf{Variables $(T, V)$ :} $dS = \frac{dU + P dV}{T} = C_v \frac{dT}{T} + nR \frac{dV}{V}$
        $$\Delta S = C_v \ln\left(\frac{T_1}{T_0}\right) + nR \ln\left(\frac{V_1}{V_0}\right)$$
        \item \textbf{Variables $(T, P)$ :} $dS = \frac{dH - V dP}{T} = C_p \frac{dT}{T} - nR \frac{dP}{P}$
        $$\Delta S = C_p \ln\left(\frac{T_1}{T_0}\right) - nR \ln\left(\frac{P_1}{P_0}\right)$$
    \end{itemize}
\end{proposition}

% \warningInfo{Chemin par étapes}{Si tu décomposes en une isochore $(0 \to A)$ puis une isobare $(A \to 1)$ comme dans ton exemple :
% $$\Delta S = C_v \ln\left(\frac{T_A}{T_0}\right) + C_p \ln\left(\frac{T_1}{T_A}\right)$$
% C'est parfaitement correct et cela redonne les formules ci-dessus en remplaçant $T_A$. }


\subsection{Démonstration du chemin par étapes}

On souhaite aller de l'état $0$ $(P_0, V_0, T_0)$ à l'état $1$ $(P_1, V_1, T_1)$. On définit un état intermédiaire $A$ tel que :
\begin{itemize}
    \item Étape 1 ($0 \to A$) : Transformation isochore ($V_A = V_0$). Pour atteindre la pression $P_1$, la température devient $T_A = \frac{P_1 V_0}{nR}$.
    \item Étape 2 ($A \to 1$) : Transformation isobare ($P_A = P_1$). On fait varier le volume de $V_0$ à $V_1$.
\end{itemize}



\subsubsection{Calcul des variations d'entropie}

1. Sur l'isochore ($0 \to A$) :
   Comme $dV = 0$, l'identité thermodynamique $dU = TdS - PdV$ devient $dU = TdS$.
   $$\Delta S_{0A} = \int_{T_0}^{T_A} \frac{C_v dT}{T} = C_v \ln\left(\frac{T_A}{T_0}\right)$$

2. Sur l'isobare ($A \to 1$) :
   Comme $dP = 0$, l'identité $dH = TdS + VdP$ devient $dH = TdS$.
   $$\Delta S_{A1} = \int_{T_A}^{T_1} \frac{C_p dT}{T} = C_p \ln\left(\frac{T_1}{T_A}\right)$$

3. Somme totale :
   $$\Delta S_{total} = C_v \ln\left(\frac{T_A}{T_0}\right) + C_p \ln\left(\frac{T_1}{T_A}\right)$$


\subsubsection{Récupération de la formule classique}

On utilise la relation de Mayer $C_p = C_v + nR$ pour décomposer le second terme :
\begin{flalign*}
    \Delta S &= C_v \ln\left(\frac{T_A}{T_0}\right) + (C_v + nR) \ln\left(\frac{T_1}{T_A}\right)\\
&= C_v \left[ \ln\left(\frac{T_A}{T_0}\right) + \ln\left(\frac{T_1}{T_A}\right) \right] + nR \ln\left(\frac{T_1}{T_A}\right)\\
&= C_v \ln\left(\frac{T_1}{T_0}\right) + nR \ln\left(\frac{T_1}{T_A}\right)
\end{flalign*}

Or, sur l'isobare $A \to 1$, on a $\frac{V}{T} = const \implies \frac{T_1}{T_A} = \frac{V_1}{V_A}$. Comme $V_A = V_0$ par construction :
$$\Delta S = C_v \ln\left(\frac{T_1}{T_0}\right) + nR \ln\left(\frac{V_1}{V_0}\right)$$

\end{document}


